{\small\opmaakproza
\begin{enumerate}[leftmargin=1.5em,itemsep=3pt,label=\textcolor{kleur1}{\sffamily\bfseries\arabic*}]
\item \textbf{Lees de inleiding en de schema's goed.} Bij een blokschema of een proces staat de informatie die je
  bij latere vragen nodig hebt vaak al in de inleiding van de opgave.
\item \textbf{Klopt je reactievergelijking?} Controleer of links en rechts evenveel atomen van elke soort staan en
  of de lading klopt. Gebruik de formules die de vraag noemt, en geef toestandsaanduidingen alleen als erom gevraagd wordt.
\item \textbf{Zoek op in je informatieboek.} Molaire massa's, oplosbaarheid, zuur- en baseconstanten, normpotentialen
  en vormingswarmtes staan in Binas of ScienceData. Noem het tabelnummer als je een waarde opzoekt.
\item \textbf{Schrijf berekeningen in stappen op}: van gegeven via mol naar het gevraagde. Elke goede stap levert
  vaak een punt op, ook als je eindantwoord niet klopt.
\item \textbf{Let op significante cijfers en eenheden.} Geef de uitkomst in het aantal cijfers dat past bij de gegevens.
  Bij een pH tellen alleen de cijfers achter de komma.
\item \textbf{Teken structuurformules volledig}, met alle atomen en bindingen die de vraag vraagt. Bij een
  \emph{gedeelte} van een polymeer laat je de uiteinden open met een streepje.
\item \textbf{Gebruik vaktaal.} Noem deeltjes bij hun naam (moleculen, ionen, atomen) en gebruik begrippen als
  waterstofbrug, polair en evenwicht precies. Een antwoord op macroniveau als het om deeltjes gaat, telt vaak niet.
\item \textbf{Bij \emph{leg uit} of \emph{verklaar}} geef je de redenering in stappen: wat gebeurt er met de deeltjes,
  en waarom leidt dat tot het verschijnsel?
\item \textbf{Kijk naar de punten.} Een vraag van 3 punten vraagt meestal drie elementen. Houd ongeveer
  2 minuten per punt aan en ga door als je vastloopt.
\item \textbf{Controleer je antwoord}: is de uitkomst logisch? Een rendement boven 100\,\% of een pH van 16 klopt niet.
\end{enumerate}\par}
